% SEGMENT OLP-0608-S01
% Part: methods
% Chapter: proofs
% Section: example-2

\documentclass[../../../include/open-logic-section]{subfiles}

\begin{document}

\olfileid{mth}{prf}{ex2}

\olsection{மற்றோர் எடுத்துக்காட்டு}

\begin{prop}
$A \subseteq C$ எனில் $A \cup (C \setminus A) = C$.
\end{prop}

% SEGMENT OLP-0608-S02
\begin{proof}
  $A \subseteq C$ என்க. $A \cup (C \setminus A) = C$
  என்று காட்ட வேண்டும்.
  \begin{quote}
    இது ஒரு நிபந்தனைக் கூற்று என்பதை முதலில் கவனிக்கிறோம்.
    மறைமுகமாக அது அனைத்தையும் பற்றிய அளவையைக் கொண்டது:
    எந்தக் கணங்கள் $A$, $C$ க்கும் இந்த முன்மொழிவு பொருந்தும்.
    எனவே $A$, $C$ தன்னிச்சையான கணங்களுக்கான மாறிகள்.
    இத்தகைய கூற்றை நிறுவ, நிபந்தனையின் முற்பகுதியை
    எடுகோளாக எடுத்துக்கொண்டு அதன் பிற்பகுதியை நிறுவ வேண்டும்.

    அடுத்து $A \subseteq C$ என்ற எடுகோளைப் பயன்படுத்துகிறோம்.
    $\subseteq$ இன் வரையறையை விரித்தால், $A$ இன்
    எல்லா !!{element}s உம் $C$ இன் !!{element}s என்று
    எடுகோள் கூறுகிறது. இதை எழுதிவைப்போம்; நிறுவல் முழுவதும்
    பயன்படும் முக்கியமான உண்மை இது.
  \end{quote}
  $\subseteq$ இன் வரையறையின்படி, $A \subseteq C$
  என்பதால், எல்லா $z$ க்கும் $z \in A$ எனில்
  $z \in C$.
  \begin{quote}
    எடுகோளில் தரப்பட்ட வரையறைகளை எல்லாம் விரித்துவிட்டோம்.
    இப்போது முடிவுக்கு வரலாம். $A \cup (C \setminus A) = C$
    என்று காட்ட வேண்டும். முந்தைய எடுத்துக்காட்டைப் போல
    நிறுவலை அமைப்போம்: $A \cup (C \setminus A)$ இன்
    ஒவ்வோர் !!{element} உம் $C$ இன் !!a{element}
    என்றும், மறுதிசையில் $C$ இன் ஒவ்வோர் !!{element}
    உம் $A \cup (C \setminus A)$ இன் !!a{element}
    என்றும் காட்டுவோம். இதைச் சுருக்கமாக
    $A \cup (C \setminus A) \subseteq C$ மற்றும்
    $C \subseteq A \cup (C \setminus A)$ என எழுதலாம்.
    (இங்கு வரையறையை விரிப்பதற்குப் பதிலாகச் சுருக்குகிறோம்;
    இது நிறுவலை வாசிக்கச் சற்றே எளிதாக்குகிறது.)
    இது இணைப்புக் கூற்று என்பதால் இரு பகுதிகளையும்
    நிறுவ வேண்டும். முதல் பகுதியைக் காட்ட, அதாவது
    $A \cup (C \setminus A)$ இன் ஒவ்வோர் !!{element}
    உம் $C$ இன் !!a{element} என்று காட்ட, தன்னிச்சையான
    $z$ க்கு $z \in A \cup (C \setminus A)$ என
    எடுத்துக்கொண்டு $z \in C$ என்று காட்டுவோம்.
    $\cup$ இன் வரையறையின்படி, $z \in A \cup (C \setminus A)$
    என்பதிலிருந்து $z \in A$ அல்லது $z \in C \setminus A$
    என்று முடிக்கலாம். இந்த முறை இப்போது பழக்கமாகத்
    தொடங்கியிருக்கும்.
  \end{quote}
  % SOURCE CORRECTION TA-MTH-003: இரண்டாம் உட்கணக் கூற்றின் விடுபட்ட மூடும் அடைப்பு சேர்க்கப்பட்டது.
  $A \cup (C \setminus A) = C$ எனில், எனில் மட்டுமே,
  $A \cup (C \setminus A) \subseteq C$ மற்றும்
  $C \subseteq (A \cup (C \setminus A))$. முதலில்
  $A \cup (C \setminus A) \subseteq C$ என்று நிறுவுவோம்.
  $z \in A \cup (C \setminus A)$ என்க. ஆகவே
  $z \in A$ அல்லது $z \in (C \setminus A)$.
  \begin{quote}
    பிரிப்புக் கூற்றை அடைந்தோம்; அதிலிருந்து
    $z \in C$ என்று நிறுவ வேண்டும். நிகழ்வுகளைப்
    பிரித்து இதைச் செய்கிறோம்.
  \end{quote}

% SEGMENT OLP-0608-S03
  நிகழ்வு 1: $z \in A$. எல்லா $z$ க்கும்
  $z \in A$ எனில் $z \in C$ என்பதால், இங்கு
  $z \in C$.
  \begin{quote}
    முன்மொழிவின் $A \subseteq C$ எடுகோளிலிருந்து
    முன்னர் பதிவுசெய்த உண்மையைப் பயன்படுத்தினோம்.
    முதல் நிகழ்வு முடிந்தது. இரண்டாவது நிகழ்வு
    $z \in (C \setminus A)$. $C \setminus A$
    என்பது இரு கணங்களின் \emph{வேறுபாடு}:
    $A$ இன் !!{element}s அல்லாத $C$ இன்
    எல்லா !!{element}s உம் உள்ள கணம். ஆனால்
    $A$ இல் இல்லாத $C$ இன் எந்த !!{element}
    உம் குறிப்பாக $C$ இன் !!a{element} ஆகும்.
  \end{quote}
  நிகழ்வு 2: $z \in (C \setminus A)$.
  இதன் பொருள் $z \in C$ மற்றும் $z \notin A$.
  எனவே குறிப்பாக $z \in C$.
  \begin{quote}
    முதல் திசையை நிறுவிவிட்டோம். இப்போது இரண்டாவது
    திசை: $C \subseteq A \cup (C \setminus A)$
    என்று நிறுவ வேண்டும். ஆகவே $z \in C$
    என எடுத்துக்கொண்டு $z \in A \cup (C \setminus A)$
    என்று காட்டுவோம்.
  \end{quote}

% SEGMENT OLP-0608-S04
  இப்போது $z \in C$ என்க. $z \in A$ அல்லது
  $z \in C \setminus A$ என்று காட்ட வேண்டும்.
  \begin{quote}
    $A$ இன் எல்லா !!{element}s உம் $C$ இன்
    !!{element}s; $C \setminus A$ என்பது
    $A$ இல் இல்லாத $C$ இன் !!{element}s கொண்ட கணம்.
    ஆகவே $z$ ஆனது $A$ இலோ $C \setminus A$
    இலோ இருக்கும். முடிவு ஏன் மெய் என்பது ஏற்கெனவே
    தெரியாவிட்டால் இந்த வாதம் தெளிவற்றதாகத் தோன்றலாம்.
    படிப்படியாக நிறுவுவது சிறந்தது. நிறுவாமல் பயன்படுத்தக்கூடிய
    எளிய உண்மை உதவும்: $z \in A$ அல்லது $z \notin A$.
    இதை ``விலக்கப்பட்ட நடுவின் கொள்கை'' என்கிறோம்:
    எந்தக் கூற்று $p$ க்கும், $p$ அல்லது அதன் மறுப்பு
    மெய். (இங்கு $p$ என்பது $z \in A$ என்ற கூற்று.)
    இது பிரிப்புக் கூற்று என்பதால் மீண்டும் நிகழ்வுகளைப்
    பிரித்து நிறுவலாம்.
  \end{quote}

% SEGMENT OLP-0608-S05
  $z \in A$ அல்லது $z \notin A$. முதல் நிகழ்வில்
  $z \in A \cup (C \setminus A)$. இரண்டாவது நிகழ்வில்
  $z \in C$ மற்றும் $z \notin A$; ஆகவே
  $z \in C \setminus A$. இதனால்
  $z \in A \cup (C \setminus A)$.
  \begin{quote}
    நிறுவல் முடிந்தது: $A \cup (C \setminus A) = C$
    என்று காட்டியுள்ளோம்.
  \end{quote}
\end{proof}

\end{document}
